% Generated draft: requires mathematical and source review.
\documentclass[11pt]{article}
\usepackage[margin=25mm]{geometry}
\usepackage{fontspec}
\setmainfont{Latin Modern Roman}
\usepackage{amsmath,amssymb,mathtools,unicode-math}
\setmathfont{Latin Modern Math}
\usepackage{enumitem,needspace,xurl,hyperref}
\setlength{\parindent}{0pt}
\setlength{\parskip}{5pt}
\setlength{\emergencystretch}{3em}
\newcommand{\N}{\mathbb{N}}
\newcommand{\Z}{\mathbb{Z}}
\newcommand{\Q}{\mathbb{Q}}
\newcommand{\R}{\mathbb{R}}
\newcommand{\C}{\mathbb{C}}
\newcommand{\F}{\mathbb{F}}
\newcommand{\A}{\mathbb{A}}
\newcommand{\PP}{\mathbb{P}}
\newcommand{\E}{\mathbb{E}}
\newcommand{\eps}{\varepsilon}
\newcommand{\dd}{\,\mathrm{d}}
\begin{document}
\section*{UnsolvedMath ID: 30002258}
\textbf{Source problem code:} catalogue entry
\subsection{Definitions and mathematical statement}
Let \(T\) be an NIP theory, let \(M\models T\) be saturated, and let \(\varphi(x,y)\) be a formula. Here NIP means that no formula \(\psi(u,v)\) has the independence property: there do not exist tuples \((a_i)_{i\in\N}\) and \((b_S)_{S\subseteq\N}\) in an elementary extension such that \(\psi(a_i,b_S)\) holds exactly when \(i\in S\). No particular saturation cardinal is specified in the supplied record.

For a parameter-free \(L(T)\)-formula \(\varphi(x,y)\), let \(S_\varphi(M)\) denote the set of complete consistent choices of instances \(\varphi(x,b)\) or \(\neg\varphi(x,b)\), for all tuples \(b\in M^{|y|}\). Equivalently, these are complete types on the Boolean algebra generated by the \(\varphi\)-instances over \(M\). The action of \(\sigma\in\operatorname{Aut}(M)\) is
\[
\sigma\cdot p
 =\{\varphi(x,\sigma(b))^{\epsilon}:
       \varphi(x,b)^{\epsilon}\in p\},
\]
where \(\theta^1=\theta\) and \(\theta^0=\neg\theta\).

Write \(\mathcal N_\varphi(M)\) for the \(\varphi\)-types over \(M\) that do not fork over \(\varnothing\), as intended in the source. One standard local interpretation is that no finite conjunction of members of \(p\) forks over \(\varnothing\), with forking understood in the usual model-theoretic sense. The supplied record does not specify whether this local interpretation or extendability to a complete nonforking type is intended; this qualification remains unresolved.

\textbf{Question.} For every such \(T\), \(M\), and \(\varphi\), does there exist a nearly Lie group \(G\) and a group homomorphism
\[
h:\operatorname{Aut}(M)\longrightarrow G
\]
such that
\[
\forall\sigma\in\operatorname{Aut}(M)\quad
h(\sigma)=1_G\ \Longrightarrow\
\forall p\in\mathcal N_\varphi(M)\quad \sigma\cdot p=p?
\]
Equivalently, the required kernel condition is
\[
\ker h\subseteq
\bigcap_{p\in\mathcal N_\varphi(M)}
\{\sigma\in\operatorname{Aut}(M):\sigma\cdot p=p\}.
\]
Thus each element of the kernel must fix each indicated type individually. Equality with this intersection is not required by the supplied statement.

\textbf{Unresolved qualifications.} The supplied record does not define ``nearly Lie group.'' This term is retained without choosing a definition, so the target class of groups must be clarified before the question is fully precise. No topology on \(\operatorname{Aut}(M)\), continuity of \(h\), or surjectivity of \(h\) is specified. The displayed action uses the parameter-free reading of \(\varphi\); if parameters in \(\varphi\) are intended, their treatment must be specified to obtain an action of the full automorphism group.
\subsection{Short English statement}
For every saturated model of an NIP theory and every formula \(\varphi\), ask whether its automorphism group admits a homomorphism into a nearly Lie group such that every element of the kernel fixes every nonforking \(\varphi\)-type over the model. The meaning of ``nearly Lie'' and the precise local nonforking convention require clarification. The supplied status is partially solved; this draft does not independently verify that assessment.
\subsection{Sources}
\begin{itemize}
\item Supplied local UnsolvedMath registry record, permanent ID \(30002258\), title ``Nearly-Lie Quotients of Automorphism Groups in NIP Theories,'' exact source problem code \texttt{catalogue entry}. The record gives proposed year 2013 and names no proposer.
\item Supplied report citation: \emph{Model Theory: Groups, Geometry, and Combinatorics} (2013), Oberwolfach Reports, pp.\ 5--66, DOI 10.4171/owr/2013/01. \url{https://doi.org/10.4171/owr/2013/01}. This is also the supplied literature reference labelled ``Model Theory, Oberwolfach Report 01/2013.'' The background describes the reports as documenting mathematical workshops and their open-problem discussions; no exact problem page is supplied.
\item Supplied literature reference: \emph{Approximate Equivalence Relations}. \url{https://ora.ox.ac.uk/objects/uuid:d079a1e5-75e9-49c2-bb91-89da02759474}. No author, date, or relevant theorem location is supplied for this reference.
\item Literature assessment dated exactly \texttt{2026-08-21}: the supplied curation labels the problem ``partially solved'' and states that the report proves nearly-Lie quotients for definable groups with generically stable measures, whereas the proposed automorphism-group homomorphism with the stated kernel condition is not established in general. This is supplied dated curation derived from the OWR record, not an independently verified literature assessment; the record expressly calls for independent verification of time-sensitive status and cited claims.
\item The registry labels the statement \texttt{corrected\_verified} and supplies the verification claim that the official report states this automorphism-group question with the nonforking type condition. That label and claim are source metadata, not independent verification in this draft.
\end{itemize}
\end{document}
